% Glossaire des plateformes (type=plateformes).
% Valeurs chiffrées : \citep[§x]{paper_id} si elles viennent du corpus (fiche ou benchmarks.csv),
% \constr si elles viennent d'une fiche constructeur (hors corpus).
% Une entrée parente est définie avant ses enfants ; profondeur maximale : 3 niveaux.

% ---------------------------------------------------------------- IP d'accélération
\newglossaryentry{dpu}{type=plateformes, name={DPU (DPUCZDX8G)}, sort={DPU},
description={IP d'inférence CNN de Xilinx (AMD) livrée avec Vitis~AI : processeur programmable
en INT8, implanté dans la PL des Zynq UltraScale+, auquel le compilateur XIR envoie un flot
d'instructions. Configurations B512 à B4096, où le nombre est le parallélisme crête en
opérations par cycle (B4096 : 4\,096 op./cycle, soit environ 1,2\,TOPS à 300\,MHz \deduit),
jusqu'à quatre cœurs par puce \constr. Sur XCZU7EV, une B4096 occupe 55,8\,k LUT, 710 DSP,
142 BRAM et 46 URAM, et la puissance estimée par Vivado va de 4,79\,W (B512) à 8,76\,W (B4096)
\citep[Table~III]{arxiv-2607.13106-no-attention-no-problem-dpu}\pp. C'est le flot
« industriel » du corpus : détection YOLO sur ZCU104 et KV260, comparaisons FPGA / Jetson,
tête d'apprentissage continu couplée au DPU \citep{2021-piyasena-accelerating-continual-learning-edge-fpga}.}}

% ---------------------------------------------------------------- SoC FPGA Xilinx (AMD)
\newglossaryentry{zynq}{type=plateformes, name={Zynq}, sort={Zynq},
description={Famille de SoC FPGA de Xilinx (AMD) : un processeur ARM durci (PS) et une logique
programmable (PL) sur la même puce, reliés par des bus AXI. Le corpus en utilise deux
générations, le Zynq-7000 et le Zynq UltraScale+ MPSoC (dont le module Kria K26). C'est la
cible principale de la thèse.}}

\newglossaryentry{zynq7000}{type=plateformes, name={Zynq-7000}, sort={Zynq-7000},
description={Première génération de Zynq, en 28\,nm
\citep[§III]{2019-libano-selective-hardening-neural-networks-fpgas} ; PS double Cortex-A9
(simple cœur pour les variantes S) et PL de type Artix-7 ou Kintex-7 selon la taille \constr.
Petite cible bas coût du corpus : FINN, concours DAC-SDC, autoencodeurs HLS, ONLAD, YOLOv2 en
RTL \citep{2025-kim-design-implementation-yolov2-accelerator-zynq}, et essais sous neutrons
(LANSCE) de \citet[§III-B]{2019-libano-selective-hardening-neural-networks-fpgas}.}}

\newglossaryentry{xc7z020}{type=plateformes, name={XC7Z020 (Zynq-7020)}, sort={XC7Z020},
parent=zynq7000,
description={Puce Zynq-7000 la plus utilisée du corpus : 53\,200 LUT, 106\,400 FF, 220 DSP
\citep[§3.1]{2025-kim-design-implementation-yolov2-accelerator-zynq}, 280 BRAM18
\citep[Table~1]{2019-tsukada-neural-network-device-learning-anomaly}, soit 140 BRAM36 et
4,9\,Mbit \constr, que \citet[§2.2]{2021-xu-dac-sdc-low-power-object} arrondissent à
630\,Ko. PS double Cortex-A9, à 650\,MHz sur PYNQ-Z1
\citep[Table~1]{2019-tsukada-neural-network-device-learning-anomaly}. Cartes du corpus :
PYNQ-Z1, PYNQ-Z2, Zybo Z7-20.}}

\newglossaryentry{pynq-z1}{type=plateformes, name={PYNQ-Z1}, sort={PYNQ-Z1}, parent=xc7z020,
description={Carte de développement Digilent \constr{} à XC7Z020-1CLG400C, 512\,Mo de DDR3,
« environ 1-4\,W » \citep[§2.2]{2021-xu-dac-sdc-low-power-object}. Carte FPGA imposée du
concours DAC-SDC ; 3,1\,W carte entière pour ONLAD
\citep[§6.4]{2019-tsukada-neural-network-device-learning-anomaly} ; 2,0 à 2,9\,W pour un
détecteur FINN \citep[§V]{arxiv-2609.24757-npu-accelerator-quantized-real-time}\pp. Utilisée
aussi par FINN-R et EF-Train.}}

\newglossaryentry{pynq-z2}{type=plateformes, name={PYNQ-Z2}, sort={PYNQ-Z2}, parent=xc7z020,
description={Carte de développement TUL à XC7Z020 (xc7z020clg400-1)
\citep[§8]{arxiv-2101.05108-fast-convolutional-neural-networks-fpgas} ; 512\,Mo de DDR3
\constr. Seule carte FPGA du corpus mesurée selon le protocole MLPerf Tiny (Joulescope), avec
hls4ml et FINN \citep[§4.4]{arxiv-2206.11791-open-source-fpga-ml-codesign}\pp ; sert aussi au
CNN QKeras 7 bits de \citet{arxiv-2101.05108-fast-convolutional-neural-networks-fpgas} et au
CAE vibratoire de \citet{2022-malviya-edge-compatible-convolutional-autoencoder-implemented}.}}

\newglossaryentry{xc7z035}{type=plateformes, name={XC7Z035}, sort={XC7Z035}, parent=zynq7000,
description={Puce Zynq-7000 de milieu de gamme : 171\,900 LUT, 343\,800 FF, 900 DSP,
500 BRAM36 \constr. Un seul usage mesuré : l'accélérateur YOLO et suivi de véhicules de
\citet[§4.1]{2023-zhai-fpga-vehicle-detection-tracking-accelerator} (XC7Z035-FFG676-2).}}

\newglossaryentry{zc706}{type=plateformes, name={ZC706}, sort={ZC706}, parent=zynq7000,
description={Carte d'évaluation Xilinx à Zynq-7000 XC7Z045 : double Cortex-A9, 218\,600 LUT,
545 BRAM \citep[§5.1]{arxiv-1612.07119-finn-framework-fast-scalable-binarized}\pp,
900 DSP \constr. Carte du premier FINN, dont la puissance murale est dominée par environ 7\,W
de repos \citep[§5.1]{arxiv-1612.07119-finn-framework-fast-scalable-binarized}\pp ; la même
puce (Z7045) porte le VAE de \citet{2024-que-low-latency-variational-autoencoder-fpgas}.}}

\newglossaryentry{zynqmp}{type=plateformes, name={Zynq UltraScale+ MPSoC}, sort={Zynq UltraScale+},
description={Seconde génération de Zynq, en 16\,nm FinFET
\citep[§III]{2019-libano-selective-hardening-neural-networks-fpgas} ; PS quadri Cortex-A53,
double Cortex-R5F et GPU Mali-400 (variantes EG et EV), encodeur vidéo VCU sur les EV ; PL à
DSP48E2, BRAM et, sur les plus grosses puces, URAM \constr. Cible du DPU de Vitis~AI. Dans le
texte, « UltraScale » seul peut aussi désigner les FPGA Kintex et Virtex UltraScale(+) sans
processeur (VCU110, VCU118, Alveo).}}

\newglossaryentry{xczu3eg}{type=plateformes, name={XCZU3EG}, sort={XCZU3EG}, parent=zynqmp,
description={Plus petite puce Zynq UltraScale+ du corpus : 70\,560 LUT, 141\,120 FF, 360 DSP,
216 BRAM36, sans URAM \constr ; « 71\,kLUT, 432 BRAM » (BRAM18) selon
\citet[§4]{arxiv-1809.04570-finn-r-end-end-deep}\pp. CodeNet en occupe la totalité des BRAM et
des DSP \citep[§5]{arxiv-2006.08357-codenet-efficient-deployment-input-adaptive}\pp. Puce de
l'Ultra96, proche de la gamme Kria.}}

\newglossaryentry{ultra96}{type=plateformes, name={Ultra96}, sort={Ultra96}, parent=xczu3eg,
description={Carte Avnet à XCZU3EG, 2\,Go de LPDDR4 \constr. Carte des éditions 2019 et
suivantes de DAC-SDC : SkyNet (pic théorique 144\,GOPS à 200\,MHz, 7,26\,W,
\citealp[§6.4]{arxiv-1909.09709-skynet-hardware-efficient-method-object}\pp), CodeNet
(4,3\,W à vide, 5,6\,W en charge, carte,
\citealp[§5]{arxiv-2006.08357-codenet-efficient-deployment-input-adaptive}\pp) et FINN-R.}}

\newglossaryentry{xczu7ev}{type=plateformes, name={XCZU7EV}, sort={XCZU7EV}, parent=zynqmp,
description={Puce Zynq UltraScale+ EV : 504\,k cellules logiques, 230\,400 LUT, 460\,800 FF,
1\,728 DSP, 312 BRAM36 et 96 URAM \constr. \citet[§IV-C]{arxiv-2301.07050-energy-efficient-reconfigurable-autoencoder-implementation}
annoncent « 504\,K LUT », ce qui est en fait le nombre de cellules logiques. Puce de la ZCU104.}}

\newglossaryentry{zcu104}{type=plateformes, name={ZCU104}, sort={ZCU104}, parent=xczu7ev,
description={Carte d'évaluation Xilinx à XCZU7EV, 2\,Go de DDR4 côté PS \constr. Carte FPGA la
plus mesurée du corpus : DPU B512 à B4096 pour les YOLO récents
\citep{arxiv-2607.13106-no-attention-no-problem-dpu}\pp, comparaison à la Jetson AGX Xavier
(23\,W à l'entrée de la carte,
\citealp[Table~1]{arxiv-2304.11580-framework-benchmarking-real-time-embedded}), RetinaNet sur
deux cœurs DPU (25\,FPS,
\citealp[§3]{arxiv-2211.11647-benchmarking-edge-computing-devices-grape}), SATAY, injection de
fautes accélérée \citep{arxiv-1912.07394-efficient-error-tolerant-quantized-neural}.}}

\newglossaryentry{xczu9eg}{type=plateformes, name={XCZU9EG}, sort={XCZU9EG}, parent=zynqmp,
description={Puce Zynq UltraScale+ EG : 274\,080 LUT, 548\,160 FF, 2\,520 DSP, 912 BRAM36,
sans URAM \constr ; 2\,520 DSP aussi dans
\citet[Table~2]{2024-yan-fpga-yolov5-accelerator-real-time}. Puce de la ZCU102, irradiée au
ChipIR par \citet[§III-B]{2019-libano-selective-hardening-neural-networks-fpgas}.}}

\newglossaryentry{zcu102}{type=plateformes, name={ZCU102}, sort={ZCU102}, parent=xczu9eg,
description={Carte d'évaluation Xilinx à XCZU9EG, 4\,Go de DDR4
\citep[Table~2]{2024-yan-fpga-yolov5-accelerator-real-time}. Sert à YOLOv5s INT4 (79\,FPS,
12,3\,W, \citealp[§7.4]{2024-yan-fpga-yolov5-accelerator-real-time}), à l'apprentissage
continu (DPU B4096 et SLDA, \citealp{2021-piyasena-accelerating-continual-learning-edge-fpga} ;
Chameleon, \citealp{2024-aggarwal-chameleon-dual-memory-replay-online} ; EF-Train) et à SMOF.}}

\newglossaryentry{xck26}{type=plateformes, name={XCK26 (Kria K26)}, sort={XCK26}, parent=zynqmp,
description={Puce Zynq UltraScale+ (variante EV, avec VCU) du module Kria K26, non vendue seule
\constr : 256\,k cellules logiques, 117\,120 LUT, 234\,240 FF, 1\,248 DSP \constr ; 144 BRAM
(648\,Ko), 64 URAM (environ 2,25\,Mo) et 4\,Go de DDR4 sur le module
\citep[Table~6]{2026-fata-low-cost-fpga-enhanced-cnn}.}}

\newglossaryentry{kv260}{type=plateformes, name={Kria KV260}, sort={KV260}, parent=xck26,
description={Kit de démarrage AMD Xilinx « Vision AI » : module K26 sur carte porteuse à
interfaces caméra. Cible matérielle visée par la thèse. Mesures du corpus : YOLOv3-Tiny INT8
sur DPU, 24,3\,FPS pour 2,13\,W (portée non précisée,
\citealp[Table~10]{2026-fata-low-cost-fpga-enhanced-cnn}) ; YOLOv5 INT8
\citep{arxiv-2507.18174-real-time-object-detection-classification}\pp ; RetinaNet sur un cœur
DPU \citep{arxiv-2211.11647-benchmarking-edge-computing-devices-grape}. Aucun protocole commun
de mesure entre ces travaux.}}

% ---------------------------------------------------------------- FPGA sans processeur durci
\newglossaryentry{artix7}{type=plateformes, name={Artix-7}, sort={Artix-7},
description={Famille FPGA Xilinx 28\,nm bas coût, sans processeur durci \constr. XC7A100T :
63\,400 LUT, 240 DSP, 135 BRAM36 \constr (135 BRAM aussi chez
\citealp[§VIII-A]{2021-manjunath-energy-efficient-edgeai-autoencoder-accelerator}) ;
XC7A35T : 20\,800 LUT, 90 DSP, 50 BRAM36 \constr. Autoencodeur de Manjunath à 100\,MHz, Arty
A7-100T sous MLPerf Tiny \citep{arxiv-2206.11791-open-source-fpga-ml-codesign}\pp, hls4ml
\citep{arxiv-2512.01463-hls4ml-flexible-open-source-platform}\pp.}}

\newglossaryentry{kintex7}{type=plateformes, name={Kintex-7}, sort={Kintex-7},
description={Famille FPGA Xilinx 28\,nm de milieu de gamme, sans processeur durci \constr.
Usage : carte KC705 de \citet{2024-zhang-end-end-acceleration-yolo-object}. « Kintex » seul
désigne aussi la Kintex UltraScale annoncée par
\citet[§6.1]{2024-campos-end-end-codesign-hessian-aware}, dont la référence (xcu250) est en
fait celle de l'Alveo U250.}}

\newglossaryentry{xc7k325t}{type=plateformes, name={XC7K325T}, sort={XC7K325T}, parent=kintex7,
description={Puce Kintex-7 : 203\,k LUT, 840 DSP, 16\,Mbit de BRAM, sur carte KC705 avec 1\,Go
de DDR3 \citep[§5.1]{2024-zhang-end-end-acceleration-yolo-object} ; 445 BRAM36 \constr. Porte
YOLOv2 et YOLOv3-Tiny en INT8 tout sur FPGA, 401 à 418\,GOP/s, 7,8\,W puce et 12,3\,W carte
\citep[§5.3]{2024-zhang-end-end-acceleration-yolo-object}.}}

\newglossaryentry{virtex7}{type=plateformes, name={Virtex-7}, sort={Virtex-7},
description={Famille FPGA Xilinx 28\,nm haut de gamme, sans processeur durci \constr. Cartes du
corpus : VC707 (XC7VX485T, 18,61\,W carte,
\citealp[§5.2]{2015-zhang-optimizing-fpga-accelerator-design-deep}), VC709 et ADM-PCIE-7V3
(XC7VX690T). « Virtex » seul désigne aussi les Virtex UltraScale(+) (VCU110, VCU118) et le
Virtex-6 ; aucune n'est une cible embarquée.}}

\newglossaryentry{xc7vx690t}{type=plateformes, name={XC7VX690T}, sort={XC7VX690T}, parent=virtex7,
description={Puce Virtex-7 : 433\,k LUT, 3\,600 DSP, 52,9\,Mbit de BRAM, deux DDR3 de 4\,Go sur
VC709 \citep[§5.1]{2024-zhang-end-end-acceleration-yolo-object} ; 1\,470 BRAM36 \constr. Sert à
ReQ-YOLO sur carte PCIe ADM-PCIE-7V3 (16\,Go de DDR3, 200\,MHz,
\citealp[§6.3.1]{2019-ding-req-yolo-resource-aware-efficient} ; 21\,W,
\citealp[Table~2]{2019-ding-req-yolo-resource-aware-efficient}) et à YOLOv2/v3 batch 2
(914\,GOP/s, \citealp[Table~3]{2024-zhang-end-end-acceleration-yolo-object}).}}

\newglossaryentry{xc6vlx240t}{type=plateformes, name={XC6VLX240T (Virtex-6)}, sort={XC6VLX240T},
description={Puce Virtex-6 (40\,nm) du kit ML605 : 241\,152 cellules logiques, 150\,720 LUT,
768 DSP48E1, 416 BRAM36 \constr ; les 768 multiplieurs sont aussi donnés par
\citet[Table~3]{2019-coutinho-deep-neural-network-hardware-implementation}. Seul usage :
l'autoencodeur empilé (SSAE) MNIST de Coutinho \citep[§V-D]{2019-coutinho-deep-neural-network-hardware-implementation}.}}

\newglossaryentry{vcu118}{type=plateformes, name={VCU118}, sort={VCU118},
description={Carte d'évaluation Xilinx à Virtex UltraScale+ XCVU9P (16\,nm) : 1\,182\,240 LUT
et 6\,840 DSP \citep[Fig.~6]{2022-sun-f-lic-fpga-learned-image} ; 2\,160 BRAM36 et 960 URAM
\constr. Grosse carte de laboratoire : compression d'image F-LIC, SATAY (YOLO, 23 à 67\,W,
\citealp[Table~III]{2023-montgomerie-corcoran-satay-streaming-architecture-toolflow-accelerating}),
SMOF, et VAE de déclenchement au LHC \citep{arxiv-2108.03986-autoencoders-fpgas-real-time-unsupervised}.}}

\newglossaryentry{polarfire}{type=plateformes, name={PolarFire}, sort={PolarFire},
description={Famille FPGA Microchip 28\,nm à configuration non volatile (flash SONOS), donc
insensible aux SEU de la mémoire de configuration \constr ; MPF300T : environ 300\,k éléments
logiques \constr. Seul usage : hls4ml sur MPF300TS et MPF100T, limité aux couches denses et
ReLU, environ 0,1\,W estimés
\citep[§5.2]{arxiv-2602.15751-enabling-low-latency-machine-learning}\pp. Intéresse la thèse
pour la tenue aux radiations.}}

\newglossaryentry{stratix}{type=plateformes, name={Stratix}, sort={Stratix},
description={Famille FPGA haut de gamme d'Intel (ex-Altera) ; Stratix~V en 28\,nm, Stratix~10
en 14\,nm \constr. Présente seulement par citation : baseline Stratix~10 GX d'EF-Train
(163\,GOPS, 20,6\,W, \citealp{arxiv-2202.10935-ef-train-enable-efficient-device}\pp) et
valeur projetée de 12\,TOP/s à 141,2\,W
\citep[Table~8]{arxiv-1806.01683-accelerating-cnn-inference-fpgas-survey}\pp.}}

% ---------------------------------------------------------------- GPU embarqués NVIDIA
\newglossaryentry{jetson}{type=plateformes, name={Jetson}, sort={Jetson},
description={Gamme NVIDIA de modules embarqués (SoC ARM + GPU CUDA, mémoire LPDDR partagée) et
de kits de développement ; logiciel JetPack et TensorRT (FP32, FP16, INT8), modes de
puissance réglables \constr. Référence GPU embarqué du corpus, mais peu de travaux la
comparent à un FPGA sur le même modèle.}}

\newglossaryentry{jetson-nano}{type=plateformes, name={Jetson Nano}, sort={Jetson Nano}, parent=jetson,
description={GPU Maxwell 128 cœurs (0,5\,TFLOPS), 4\,Go de LPDDR4, 5 à 10\,W
\citep[Table~2]{2022-feng-benchmark-analysis-yolo-performance-edge} ; CPU quadri Cortex-A57
\constr. Ne supporte que le FP16 selon
\citet[§3.3.1]{2022-feng-benchmark-analysis-yolo-performance-edge}. Mesures YOLO
\citep{arxiv-2409.16808-comprehensive-evaluation-deep-learning-object}\pp, apprentissage
continu \citep{2024-aggarwal-chameleon-dual-memory-replay-online}, et point de comparaison
repris par plusieurs papiers FPGA.}}

\newglossaryentry{tx2}{type=plateformes, name={Jetson TX2}, sort={Jetson TX2}, parent=jetson,
description={GPU Pascal 256 cœurs, CPU double Denver~2 et quadri Cortex-A57, 8\,Go de LPDDR4,
7,5 à 15\,W \constr ; 665\,GFLOPS crête à 1\,300\,MHz selon
\citet[§6.4]{arxiv-1909.09709-skynet-hardware-efficient-method-object}\pp. Plateforme GPU
imposée de DAC-SDC \citep[§2.2]{2021-xu-dac-sdc-low-power-object} ; baseline de SkyNet, SATAY
et ReQ-YOLO.}}

\newglossaryentry{xavier-nx}{type=plateformes, name={Jetson Xavier NX}, sort={Jetson Xavier NX}, parent=jetson,
description={GPU Volta 384 cœurs + 48 Tensor Cores (1,3\,TFLOPS), 8\,Go de LPDDR4x, 10 à 15\,W
\citep[Table~2]{2022-feng-benchmark-analysis-yolo-performance-edge} ; CPU 6 cœurs Carmel,
21\,TOPS INT8 \constr. Baseline GPU de \citet{2021-piyasena-accelerating-continual-learning-edge-fpga},
de F-LIC \citep{2022-sun-f-lic-fpga-learned-image} et de
\citet{arxiv-2407.17647-energy-efficient-artefact-detection-accelerator}\pp{} (mode 10\,W).}}

\newglossaryentry{agx-xavier}{type=plateformes, name={Jetson AGX Xavier}, sort={Jetson AGX Xavier}, parent=jetson,
description={CPU 8 cœurs Carmel, GPU et NVDLA
\citep[§V-A]{arxiv-2406.12362-certified-ml-object-detection-surveillance}\pp ; GPU Volta
512 cœurs + 64 Tensor Cores, 2 NVDLA, 32\,Go, 10 à 30\,W, 32\,TOPS INT8 \constr. 16,5 à 28\,W
mesurés à l'entrée de la carte selon la précision, face à une ZCU104
\citep[Table~1]{arxiv-2304.11580-framework-benchmarking-real-time-embedded} ; cible de la
détection certifiée en vidéosurveillance \citep{arxiv-2406.12362-certified-ml-object-detection-surveillance}\pp.}}

\newglossaryentry{orin}{type=plateformes, name={Jetson Orin}, sort={Jetson Orin}, parent=jetson,
description={Génération Jetson à GPU Ampere et CPU Cortex-A78AE \constr : Orin Nano, Orin NX,
AGX Orin. Les lignes Orin du corpus viennent de deux papiers seulement, et aucun travail ne
compare une FPGA à une Orin.}}

\newglossaryentry{orin-nano}{type=plateformes, name={Jetson Orin Nano}, sort={Jetson Orin Nano}, parent=orin,
description={6 × Cortex-A78AE, GPU Ampere 1\,024 cœurs + 32 Tensor Cores, 8\,Go, 7 à 15\,W
\citep[Table~2]{arxiv-2502.15737-performance-analysis-you-only-look} ; 40\,TOPS INT8 annoncés
au lancement pour la version 8\,Go \constr. Version 4\,Go dans
\citet[Table~I]{arxiv-2409.16808-comprehensive-evaluation-deep-learning-object}\pp, qui y
mesurent l'énergie au repos la plus élevée de leur banc (362\,mWh sur 5\,min, Fig.~2a).}}

\newglossaryentry{orin-nx}{type=plateformes, name={Jetson Orin NX}, sort={Jetson Orin NX}, parent=orin,
description={8 × Cortex-A78AE, 16\,Go, 10 à 25\,W
\citep[Table~2]{arxiv-2502.15737-performance-analysis-you-only-look}. La même table donne un GPU
Ampere de 2\,048 cœurs + 64 Tensor Cores, alors que la fiche constructeur annonce 1\,024 cœurs
+ 32 Tensor Cores et 100\,TOPS INT8 pour la version 16\,Go \constr. YOLOv8n INT8 : 65,8\,FPS
pour environ 10,3\,W \citep[Table~4]{arxiv-2502.15737-performance-analysis-you-only-look}.}}

\newglossaryentry{agx-orin}{type=plateformes, name={Jetson AGX Orin}, sort={Jetson AGX Orin}, parent=orin,
description={15 à 40\,W et 275\,TOPS selon
\citet[Table~4]{arxiv-2411.01431-efficient-deep-learning-infrastructures-embedded} ; le
constructeur annonce 12 × Cortex-A78AE et un GPU Ampere 2\,048 cœurs + 64 Tensor Cores, avec
275\,TOPS et 15 à 60\,W pour la version 64\,Go, 200\,TOPS et 15 à 40\,W pour la version 32\,Go
\constr : la table du corpus mêle les deux. Baseline GPU d'ImageHD
\citep{arxiv-2604.21280-imagehd-energy-efficient-device-continual}\pp.}}

% ---------------------------------------------------------------- MCU et SoC ultra-basse consommation
\newglossaryentry{pulp}{type=plateformes, name={PULP}, sort={PULP},
description={Parallel Ultra-Low Power : plateforme matérielle open source de l'ETH Zürich et de
l'université de Bologne, à cluster de cœurs RISC-V partageant une mémoire L1 \constr. Base des
puces VEGA et GAP9 du corpus.}}

\newglossaryentry{vega}{type=plateformes, name={VEGA}, sort={VEGA}, parent=pulp,
description={SoC PULP en 22\,nm GlobalFoundries : cluster de 9 cœurs RISC-V (8 de calcul + 1
contrôleur), 4 FPU FP32/FP16 partagées, 128\,Ko de L1, 1,5\,Mo de L2, 4\,Mo de MRAM
\citep[§IV-A]{arxiv-2110.10486-tinyml-platform-device-continual-learning}\pp ; 375\,MHz et
62\,mW mesurés sur la puce seule \citep[§V-D]{arxiv-2110.10486-tinyml-platform-device-continual-learning}\pp.
Plateforme de l'apprentissage continu par latent replay sur MCU.}}

\newglossaryentry{gap9}{type=plateformes, name={GAP9}, sort={GAP9}, parent=pulp,
description={MCU RISC-V de GreenWaves, dérivé de PULP, avec accélérateur NE16 \constr. Fabric
controller + cluster de 9 cœurs, L3 externe de 8\,Mo, 650\,mV et 240\,MHz, environ 44 à 50\,mW
\citep[§V-A, §VI-C]{arxiv-2403.07851-12-mj-per-class-device}\pp ; cluster 8 cœurs à 175\,MHz
et 1,6\,Mo de SRAM selon
\citet[Table~4]{2025-el-zeinaty-designing-object-detection-models-tinyml}. Apprentissage
few-shot à 12\,mJ par classe et TinyissimoYOLO.}}

\newglossaryentry{stm32}{type=plateformes, name={STM32}, sort={STM32},
description={Famille de MCU 32 bits de STMicroelectronics à cœur Arm Cortex-M ; déploiement de
réseaux par X-CUBE-AI ou TensorFlow Lite Micro \constr. Cible des travaux MCUNet,
TinyissimoYOLO et de la comparaison à VEGA ; la référence « petite cible » face à laquelle un
FPGA doit se justifier.}}

\newglossaryentry{stm32f412}{type=plateformes, name={STM32F412}, sort={STM32F412}, parent=stm32,
description={Cortex-M4, 256\,Ko de SRAM, 1\,Mo de Flash
\citep[§4]{arxiv-2110.15352-mcunetv2-memory-efficient-patch-inference}\pp ; 100\,MHz \constr.
Porte MCUNetV2-YOLOv3-M4 (64,6\,\% de mAP VOC,
\citealp[Table~4]{2025-el-zeinaty-designing-object-detection-models-tinyml}).}}

\newglossaryentry{stm32f746}{type=plateformes, name={STM32F746}, sort={STM32F746}, parent=stm32,
description={Cortex-M7 à 216\,MHz, 320\,Ko de SRAM, 1\,Mo de Flash
\citep[§3.1]{arxiv-2206.15472-device-training-256kb-memory}\pp{} ;
\citep[Table~4]{2025-el-zeinaty-designing-object-detection-models-tinyml}. Cible de
l'entraînement sur l'appareil en 256\,Ko et de MCUNetV2.}}

\newglossaryentry{stm32h743}{type=plateformes, name={STM32H743}, sort={STM32H743}, parent=stm32,
description={Cortex-M7 à 480\,MHz, 1\,Mo de SRAM
\citep[Table~4]{2025-el-zeinaty-designing-object-detection-models-tinyml} ; 512\,Ko de SRAM et
2\,Mo de Flash selon \citet[§4]{arxiv-2110.15352-mcunetv2-memory-efficient-patch-inference}\pp{}
(le constructeur annonce 1\,Mo de RAM au total, dont 512\,Ko de SRAM AXI d'un seul tenant
\constr). MCU de référence des détecteurs MCUNet, EtinyNet et XiNet.}}

\newglossaryentry{stm32h7a3}{type=plateformes, name={STM32H7A3}, sort={STM32H7A3}, parent=stm32,
description={Cortex-M7, utilisé à 3,3\,V et 192\,MHz
\citep[Table~II]{arxiv-2306.00001-tinyissimoyolo-quantized-low-memory-footprint}\pp ; jusqu'à
280\,MHz, 1,4\,Mo de SRAM et 2\,Mo de Flash \constr. TinyissimoYOLO y tourne en 359\,ms par
inférence \citep[§III]{arxiv-2306.00001-tinyissimoyolo-quantized-low-memory-footprint}\pp.}}

\newglossaryentry{stm32l4}{type=plateformes, name={STM32L4}, sort={STM32L4}, parent=stm32,
description={Série basse consommation à Cortex-M4. STM32L476RG à 80\,MHz (carte NUCLEO-64),
65 fois plus lent que VEGA en moyenne
\citep[§V-E]{arxiv-2110.10486-tinyml-platform-device-continual-learning}\pp ; 128\,Ko de SRAM
\constr. STM32L4R9 (Cortex-M4F, 3,3\,V, 120\,MHz) dans
\citet[Table~II]{arxiv-2306.00001-tinyissimoyolo-quantized-low-memory-footprint}\pp.}}

\newglossaryentry{stm32n6}{type=plateformes, name={STM32N6}, sort={STM32N6}, parent=stm32,
description={MCU à Cortex-M55 et NPU, 800\,MHz, 4,2\,Mo de SRAM
\citep[Table~4]{2025-el-zeinaty-designing-object-detection-models-tinyml} ; NPU Neural-ART
d'environ 600\,GOPS \constr. Cité comme MCU à accélérateur IA ; aucune mesure dans le corpus.}}

% ---------------------------------------------------------------- cœurs Arm
\newglossaryentry{cortex-m4}{type=plateformes, name={Cortex-M4}, sort={Cortex-M4},
description={Cœur Arm ARMv7E-M pour MCU, avec instructions DSP/SIMD et FPU simple précision
optionnelle (Cortex-M4F) \constr. Présent dans les STM32F412 et STM32L4, le MAX78000 et
l'Apollo4b du corpus.}}

\newglossaryentry{cortex-m7}{type=plateformes, name={Cortex-M7}, sort={Cortex-M7},
description={Cœur Arm ARMv7E-M haut de gamme pour MCU : pipeline superscalaire, caches, FPU
optionnelle \constr. Présent dans les STM32F746, STM32H743 et STM32H7A3 du corpus.}}

\newglossaryentry{cortex-a9}{type=plateformes, name={Cortex-A9}, sort={Cortex-A9},
description={Cœur applicatif Arm ARMv7-A \constr ; processeur (PS) des Zynq-7000, à 650\,MHz
sur PYNQ-Z1 \citep[Table~1]{2019-tsukada-neural-network-device-learning-anomaly}. Exécute le
post-traitement YOLOv2 sous PetaLinux chez
\citet[§4.2]{2025-kim-design-implementation-yolov2-accelerator-zynq}.}}

\newglossaryentry{cortex-a15}{type=plateformes, name={Cortex-A15}, sort={Cortex-A15},
description={Cœur applicatif Arm ARMv7-A \constr. Cible de mesure (SoC TI KeyStone) du code C
avionique généré par ACETONE \citep[§4.2]{2022-de-albuquerque-silva-acetone-predictable-programming-framework-ml}.}}

% ---------------------------------------------------------------- NPU, VPU, ASIC
\newglossaryentry{ncs2}{type=plateformes, name={Intel NCS2 (Myriad)}, sort={NCS2},
description={Neural Compute Stick 2 : clé USB~3.0 d'Intel à VPU Movidius Myriad X, environ
1\,TFLOPS et 1\,W en plus de l'hôte \citep[Table~2]{2022-feng-benchmark-analysis-yolo-performance-edge},
0,64\,W au repos \citep[Table~4]{2022-feng-benchmark-analysis-yolo-performance-edge} ;
OpenVINO, FP16. Son prédécesseur, la Myriad~2, est la VPU des satellites Φ-Sat
\citep{2023-guerrisi-artificial-intelligence-board-image-compression} et a été testée sous
ions lourds \citep[§IV-C.3]{2024-rech-artificial-neural-networks-space-safety}.}}

\newglossaryentry{hailo8}{type=plateformes, name={Hailo-8}, sort={Hailo-8},
description={NPU dataflow de la société Hailo, 26\,TOPS INT8 pour environ 2,5\,W typiques
\constr ; utilisée via la carte Raspberry Pi AI HAT+ et HailoRT
\citep[§III-B2]{arxiv-2409.16808-comprehensive-evaluation-deep-learning-object}\pp. Seul point
NPU du corpus : YOLOv8n, mAP COCO 32 en 76\,ms
\citep[§IV]{arxiv-2409.16808-comprehensive-evaluation-deep-learning-object}\pp.}}

\newglossaryentry{edgetpu}{type=plateformes, name={Google Edge TPU (Coral)}, sort={Edge TPU},
description={ASIC d'inférence INT8 de Google, 4\,TOPS pour 2\,W
\citep[Table~4]{arxiv-2411.01431-efficient-deep-learning-infrastructures-embedded}, vendu en
clé USB, module ou Coral Dev Board \constr. Mesuré sur Coral Dev Board
\citep{arxiv-2211.11647-benchmarking-edge-computing-devices-grape} et en clé USB sur Raspberry
Pi, où YOLOv8 doit passer en 320 × 320
\citep[§III-B2]{arxiv-2409.16808-comprehensive-evaluation-deep-learning-object}\pp.}}

\newglossaryentry{eyeriss}{type=plateformes, name={Eyeriss}, sort={Eyeriss},
description={Accélérateur CNN ASIC du MIT (65\,nm, 168 PE, flot « row stationary») \constr ;
166\,GOPS/W pour 278\,mW sur AlexNet \citep[§XI-B]{2023-shuvo-efficient-acceleration-deep-learning-inference}.
Architecture de référence de l'étude de propagation d'erreurs et de FIT de
\citet{2017-li-understanding-error-propagation-deep-learning}, en simulation.}}
